% Generated from locale/tr/content/normal-modal-logic/syntax-and-semantics/substitution.tex; do not hand-edit.


\olfileid[tr]{nml}{syn}{sub}

\olsection{Eşzamanlı Yerine Koyma}

!!a{formula}~$!A$'nın bir \emph{örneği}, $!A$ içindeki
!!a{propositional variable}{}nin bütün geçişlerinin başka bir !!{formula} ile
değiştirilmesiyle elde edilir. !!{formula}s{}in örneklerinden gerek geçerliliği
gerek !!{derivability} kavramını tartışırken sık sık söz edeceğiz.
Dolayısıyla bu kavramı kesin olarak tanımlamak yararlıdır.

\begin{defn}\ollabel{def:subst-inst}
  $!A$ bir kiplik !!{formula} olsun ve bütün !!{propositional variable}s
  $p_1$, \dots, $p_n$ arasında bulunsun. $!D_1$, \dots, $!D_n$ de kiplik
  !!{formula}s ise
  $\SSubst{!A}{\subst{!D_1}{p_1}, \dots, \subst{!D_n}{p_n}}$ ifadesini,
  $!A$ içinde her $p_i$ yerine $!D_i$'nin eşzamanlı olarak konulmasının
  sonucu diye tanımlarız. Biçimsel olarak bu, $!A$ üzerinde tümevarımlı bir
  tanımdır:
  \begin{enumerate}
    \tagitem{prvFalse}{\indcase{!A}{\lfalse}{%
        $\SSubst{\indfrm}{\subst{!D_1}{p_1}, \dots,
          \subst{!D_n}{p_n}}$, $\lfalse$'dır.}}{}
    \tagitem{prvTrue}{\indcase{!A}{\ltrue}{%
        $\SSubst{\indfrm}{\subst{!D_1}{p_1}, \dots,
          \subst{!D_n}{p_n}}$, $\ltrue$'dur.}}{}
    \item \indcase{!A}{q}{$i = 1$, \dots,~$n$ için $q \not\ident p_i$
      olmak koşuluyla $\SSubst{\indfrm}{\subst{!D_1}{p_1}, \dots,
        \subst{!D_n}{p_n}}$, $q$'dur.}
    \item \indcase{!A}{p_i}{$\SSubst{\indfrm}{\subst{!D_1}{p_1},
        \dots, \subst{!D_n}{p_n}}$, $!D_i$'dir.}
    \tagitem{prvNot}{\indcase{!A}{\lnot !B}{%
        $\SSubst{\indfrm}{\subst{!D_1}{p_1}, \dots, \subst{!D_n}{p_n}}$,
        $\lnot \SSubst{!B}{\subst{!D_1}{p_1}, \dots, \subst{!D_n}{p_n}}$'dır.}}{}
    \tagitem{prvAnd}{\indcase{!A}{(!B \land
        !C)}{$\SSubst{\indfrm}{\subst{!D_1}{p_1}, \dots,
          \subst{!D_n}{p_n}}$,
        \[(\SSubst{!B}{\subst{!D_1}{p_1}, \dots, \subst{!D_n}{p_n}} \land
          \SSubst{!C}{\subst{!D_1}{p_1}, \dots, \subst{!D_n}{p_n}})\]
        ifadesidir.}}{}
    \tagitem{prvOr}{\indcase{!A}{(!B \lor
          !C)}{$\SSubst{\indfrm}{\subst{!D_1}{p_1}, \dots,
          \subst{!D_n}{p_n}}$,
        \[(\SSubst{!B}{\subst{!D_1}{p_1}, \dots, \subst{!D_n}{p_n}} \lor
          \SSubst{!C}{\subst{!D_1}{p_1}, \dots, \subst{!D_n}{p_n}})\]
        ifadesidir.}}{}
    \tagitem{prvIf}{\indcase{!A}{(!B \lif
          !C)}{$\SSubst{\indfrm}{\subst{!D_1}{p_1}, \dots,
          \subst{!D_n}{p_n}}$,
        \[(\SSubst{!B}{\subst{!D_1}{p_1}, \dots, \subst{!D_n}{p_n}} \lif
          \SSubst{!C}{\subst{!D_1}{p_1}, \dots, \subst{!D_n}{p_n}})\]
        ifadesidir.}}{}
    \tagitem{prvIff}{\indcase{!A}{(!B \liff
          !C)}{$\SSubst{\indfrm}{\subst{!D_1}{p_1}, \dots,
          \subst{!D_n}{p_n}}$,
        \[(\SSubst{!B}{\subst{!D_1}{p_1}, \dots, \subst{!D_n}{p_n}} \liff
          \SSubst{!C}{\subst{!D_1}{p_1}, \dots, \subst{!D_n}{p_n}})\]
        ifadesidir.}}{}
    \tagitem{prvBox}{\indcase{!A}{\Box !B}{%
        $\SSubst{\indfrm}{\subst{!D_1}{p_1}, \dots, \subst{!D_n}{p_n}}$,
        $\Box
        \SSubst{!B}{\subst{!D_1}{p_1}, \dots, \subst{!D_n}{p_n}}$'dır.}}{}
    \tagitem{prvDiamond}{\indcase{!A}{\Diamond !B}{%
        $\SSubst{\indfrm}{\subst{!D_1}{p_1}, \dots, \subst{!D_n}{p_n}}$,
        $\Diamond
        \SSubst{!B}{\subst{!D_1}{p_1}, \dots, \subst{!D_n}{p_n}}$'dır.}}{}
  \end{enumerate}
  $\SSubst{!A}{\subst{!D_1}{p_1}, \dots,
    \subst{!D_n}{p_n}}$ ile verilen !!{formula}, $!A$'nın \emph{yerine
    koyma örneği} olarak adlandırılır.
\end{defn}

\begin{ex}
  $!A$'nın $p_1 \lif \Box(p_1 \land p_2)$, $!D_1$'in
  $\Diamond(p_2 \lif p_3)$ ve $!D_2$'nin $\lnot\Box p_1$ olduğunu
  varsayın. O zaman
  $\SSubst{!A}{\subst{!D_1}{p_1}, \subst{!D_2}{p_2}}$ şöyledir:
  \begin{align*}
    \Diamond(p_2 \lif p_3) &
    \lif \Box(\Diamond(p_2 \lif p_3) \land \lnot\Box p_1)
    \intertext{oysa $\SSubst{!A}{\subst{!D_2}{p_1}, \subst{!D_1}{p_2}}$ şöyledir:}
    \lnot\Box p_1 &
    \lif \Box(\lnot\Box p_1 \land \Diamond(p_2 \lif p_3))
    \intertext{Eşzamanlı yerine koymanın genellikle yinelenmiş yerine koymayla
      aynı olmadığına dikkat edin. Örneğin yukarıdaki
      $\SSubst{!A}{\subst{!D_1}{p_1}, \subst{!D_2}{p_2}}$ ile
      $\Subst{(\Subst{!A}{!D_1}{p_1})}{!D_2}{p_2}$ ifadesini karşılaştırın;
      ikincisi şudur:}
    &\Subst{\Diamond(p_2 \lif p_3) \lif
      \Box(\Diamond(p_2 \lif p_3) \land p_2)}{\lnot\Box p_1}{p_2},
      \text{ yani}\\
    \Diamond(\lnot\Box p_1 \lif p_3) &
    \lif \Box(\Diamond(\lnot\Box p_1
    \lif p_3) \land \lnot\Box p_1)\\
    \intertext{ve $\Subst{(\Subst{!A}{!D_2}{p_2})}{!D_1}{p_1}$ ile karşılaştırın:}
    &\Subst{p_1 \lif \Box(p_1 \land \lnot\Box p_1)}
      {\Diamond(p_2 \lif p_3)}{p_1}, \text{ yani}\\
    \Diamond(p_2 \lif p_3) &
    \lif \Box(\Diamond(p_2
    \lif p_3) \land \lnot\Box\Diamond(p_2 \lif p_3)).
  \end{align*}
\end{ex}

